% fig1_framework.tex — 顶刊级框架图（TikZ 矢量，语义化配色 + 图标 + 协议条）
% 结构对应 Methodology：数据源 -> 特征 -> 模型 -> 风控 -> 决策
\begin{figure*}[tbp]
\centering
\resizebox{\textwidth}{!}{%
\begin{tikzpicture}[
  font=\scriptsize,
  % ===== 配色（语义化，灰度友好）=====
  cDataBG/.style={fill=blue!7},   cDataBD/.style={draw=blue!55!black},
  cFeatBG/.style={fill=violet!6}, cFeatBD/.style={draw=violet!55!black},
  cModBG/.style={fill=teal!6},    cModBD/.style={draw=teal!55!black},
  cRiskBG/.style={fill=orange!8}, cRiskBD/.style={draw=orange!70!black},
  cDecBG/.style={fill=gray!8},    cDecBD/.style={draw=gray!55!black},
  % ===== 样式 =====
  dsrc/.style={draw, rounded corners=2pt, fill=white, line width=0.7pt,
               minimum height=9.5mm, minimum width=30mm, align=center},
  sub/.style={draw, rounded corners=2pt, fill=white, line width=0.7pt,
              minimum height=8mm, minimum width=27mm, align=center},
  panel/.style={rounded corners=6pt, line width=1pt, inner sep=7pt},
  stageLab/.style={font=\small\bfseries},
  arr/.style={-{Stealth[length=2.6mm, width=2.2mm]}, line width=0.9pt},
  chip/.style={rounded corners=2pt, font=\scriptsize\bfseries, text=white,
               minimum height=5.5mm, minimum width=15mm, align=center},
  note/.style={font=\scriptsize\itshape, gray!60!black},
  lab/.style={font=\scriptsize\bfseries}
]
% ===== 数据源图标（手绘迷你图）=====
\newcommand{\icondb}{\tikz[baseline=-0.4ex, scale=0.42]{%
\draw[line width=0.9pt] (0,0) ellipse (0.9 and 0.3);
\draw[line width=0.9pt] (0,0.6) ellipse (0.9 and 0.3);
\draw[line width=0.9pt] (0,1.2) ellipse (0.9 and 0.3);
\draw[line width=0.9pt] (-0.9,0.6) -- (-0.9,1.2);
\draw[line width=0.9pt] (0.9,0.6) -- (0.9,1.2);
\draw[line width=0.9pt] (-0.9,0) -- (-0.9,0.6);
\draw[line width=0.9pt] (0.9,0) -- (0.9,0.6);}}
\newcommand{\iconline}{\tikz[baseline=-0.4ex, scale=0.42]{%
\draw[line width=0.9pt] (0,0.4) -- (0.5,1.1) -- (1.0,0.6) -- (1.5,1.5) -- (2.0,0.8);
\draw[line width=0.6pt, dashed] (0,0) -- (2.2,0);}}
\newcommand{\iconbars}{\tikz[baseline=-0.4ex, scale=0.42]{%
\draw[line width=0.9pt] (0,0) rectangle (0.5,1.0);
\draw[line width=0.9pt] (0.7,0) rectangle (1.2,1.7);
\draw[line width=0.9pt] (1.4,0) rectangle (1.9,0.7);}}
\newcommand{\iconcards}{\tikz[baseline=-0.4ex, scale=0.42]{%
\draw[line width=0.9pt, rounded corners=0.6pt] (0,0) rectangle (1.0,1.4);
\draw[line width=0.9pt, rounded corners=0.6pt] (0.4,0.3) rectangle (1.4,1.7);}}

% ================= 阶段1：数据源 =================
\node[dsrc] (s1) at (0, 1.30) {{\icondb}\\ Structured\\ statistics};
\node[dsrc] (s2) at (0, 0.35) {{\iconline}\\ xG time series};
\node[dsrc] (s3) at (0, -0.60) {{\iconbars}\\ Market odds (8 books)};
\node[dsrc] (s4) at (0, -1.55) {{\iconcards}\\ Referee history};
\node[panel, cDataBG, cDataBD, fit=(s1)(s2)(s3)(s4)] (A) {};
\node[stageLab, cDataBD] at ([yshift=1.6mm]A.north) {Pre-match data};

% ================= 阶段2：特征工程 =================
\node[dsrc] (f1) at (4.9, 0.55) {\textbf{Feature card}\\ \scriptsize 60 features\\[1pt]
  \scriptsize rolling form, rank\\ \scriptsize de-vig probabilities\\ \scriptsize odds movement, xG\\ \scriptsize referee stats};
\node[chip, fill=red!70!black, minimum width=22mm] (f2) at (4.9, -1.15) {no leakage};
\node[note] at (4.9, -1.75) {information cutoff = kickoff};
\node[panel, cFeatBG, cFeatBD, fit=(f1)(f2)] (B) {};
\node[stageLab, cFeatBD] at ([yshift=1.6mm]B.north) {Feature engineering};

% ================= 阶段3：模型 =================
\node[sub] (m1) at (9.0, 1.10) {\scriptsize\textbf{Statistical}\\ Poisson, Dixon--Coles, Elo};
\node[sub] (m2) at (9.0, 0.10) {\scriptsize\textbf{Machine learning}\\ Random Forest, XGBoost};
\node[sub] (m3) at (9.0, -0.90) {\scriptsize\textbf{LLM (DeepSeek)}\\ domain-enhanced, no retrieval};
\node[panel, cModBG, cModBD, fit=(m1)(m2)(m3)] (C) {};
\node[stageLab, cModBD] at ([yshift=1.6mm]C.north) {Forecast models};

% ================= 阶段4：风控 =================
\node[dsrc] (r1) at (13.4, 1.10) {\textbf{SCS}\\ complexity score};
\node[dsrc] (r2) at (13.4, 0.10) {\textbf{UI}\\ uncertainty index};
\node[chip, fill=green!60!black] (r3) at (12.2, -1.05) {full stake};
\node[chip, fill=orange!85!black] (r4) at (13.4, -1.05) {scaled};
\node[chip, fill=red!75!black] (r5) at (14.6, -1.05) {no-bet};
\node[note] at (13.4, -1.70) {thresholds tuned on validation};
\node[panel, cRiskBG, cRiskBD, fit=(r1)(r2)(r3)(r4)(r5)] (D) {};
\node[stageLab, cRiskBD] at ([yshift=1.6mm]D.north) {Risk control};

% ================= 阶段5：决策 =================
\node[dsrc] (o1) at (16.6, 0) {\textbf{Bet decision}\\[1pt]
  \scriptsize stake or no-bet\\ \scriptsize stop-loss rule};
\node[panel, cDecBG, cDecBD, fit=(o1)] (E) {};
\node[stageLab, cDecBD] at ([yshift=1.6mm]E.north) {Decision};

% ================= 评估协议条 =================
\node[panel, cDecBG, cDecBD, minimum width=150mm] (P) at (8.3, -3.3) {
  \scriptsize\textbf{Evaluation protocol:} chronological split
  (train 2019/20--2023/24 \textbullet{} val 2024/25 \textbullet{} test 2025/26) \quad
  bootstrap 95\% CI \quad equal-stake betting \quad no-bet reported separately};

% ================= 箭头 =================
\foreach \s in {s1, s2, s3, s4} {
  \draw[arr] (\s.east) -- (f1.west);
}
\draw[arr] (f1.east) -- (m2.west);
\node[note, yshift=2mm] at ($(f1.east)!0.5!(m2.west)$) {probabilities};
\draw[arr] (m2.east) -- (r2.west);
\draw[arr] (r2.east) -- (o1.west);
\node[note, yshift=2mm] at ($(r2.east)!0.5!(o1.west)$) {position limits};
% 决策到协议条（虚线，表示评估闭环）
\draw[line width=0.7pt, dashed, gray] (E.south) -- (P.north);
\end{tikzpicture}%
}
\caption{Overview of the proposed framework. All features are constructed
from pre-match information only; every model is evaluated against de-vigged
closing odds under an identical protocol; and the risk-control layer converts
model uncertainty into stake sizing and no-bet decisions.}
\label{fig:framework}
\end{figure*}
